\subsection{Homology localizers and their exact action criterion}

\begin{proposition}[Proved: the homology localizer]\label{D:homology-localizer}
For every nonzero commutative ring $R$ with identity, let $W_R$ be the functors whose nerves induce isomorphisms on $H_j(-;R)$ for every $j\geq0$. Then $W_R$ is a fundamental localizer with the precise relative Theorem~A axiom fixed in Part~A.
\end{proposition}
\begin{proof}
Identities, two-out-of-three, and closure under retracts follow from the corresponding properties of isomorphisms of $R$-modules. A category with a terminal object has a contractible nerve and therefore maps to the point by a member of $W_R$.

For the relative axiom, consider functors $u:A\to B$, $v:B\to C$, and $f=vu:A\to C$, and assume that $A/c\to B/c$ belongs to $W_R$ for every $c\in C$. Use the bisimplicial set
\[
 S_A(p,q)=\{a_0\to\cdots\to a_q,\quad
                  f(a_q)\to c_0\to\cdots\to c_p\}.
\]
The last $A$-face composes the last arrow of the $A$-chain with the displayed arrow to $c_0$; the first $C$-face composes that displayed arrow with $c_0\to c_1$. The other faces and degeneracies have the evident nerve formulas. These formulas commute and define the indicated bisimplicial set. Define $S_B$ similarly using $v$. The functor $u$ induces a bisimplicial map $S_A\to S_B$ by applying $u$ to the $A$-chain and leaving the $C$-chain unchanged.

For a fixed $A$-chain, the $p$-direction is the nerve of $f(a_q)\backslash C$, which has initial object $\id_{f(a_q)}$. Thus projection onto the $A$-chain induces a homology isomorphism from the total normalized chains of $S_A$ to $C_*(NA;R)$. This follows by taking homology in the $p$-direction: only its degree-zero row remains. The same argument applies to $S_B$ and $B$.

For a fixed $C$-chain, the $q$-direction is the nerve of $A/c_0$. Indeed, its last arrow to $c_0$ determines the arrows from all preceding objects by composition. The analogous direction of $S_B$ is $B/c_0$. The hypothesis therefore makes $S_A\to S_B$ an isomorphism after taking homology in the $q$-direction, hence on the homology of the total complexes. Both spectral sequences are first-quadrant and converge degreewise. The two projections commute with the specified map $Nu$, so $Bu$ induces the resulting homology isomorphism. Hence $u\in W_R$.
\end{proof}

\begin{theorem}[Proved: total $R$-asphericity]\label{D:R-action}
A small category $A$ is totally $W_R$-aspherical if and only if it is nonempty and connected and satisfies $\mathcal P_R(A)$ of Definition~\ref{C:action-definition}. In particular
\[
 \{A:A\text{ is totally }W_{\mathbb Z}\text{-aspherical}\}
 =\{A:A\text{ is totally }\Winfty\text{-aspherical}\}.
\]
\end{theorem}
\begin{proof}
If $A$ is total, its identity functor first gives $R$-acyclicity. Since $R\ne0$, the induced map $R[\pi_0(A)]\to R$ is an isomorphism only when $A$ is nonempty and connected. Apply the collage construction to each component subfunctor of an arbitrary $H:A\to\Set$. The reduced Mayer--Vietoris sequence gives acyclicity of that component, exactly as in Theorem~\ref{C:classification}, with $R$ replacing $\mathbb Z$.

Conversely, for a final functor $f:A\to B$, every $b\backslash A$ is a connected category of elements and is therefore $R$-acyclic. Lemma~\ref{C:homology-A} makes $Bf$ an $R$-homology isomorphism. The constant singleton functor gives $R$-acyclicity of $A$, so $B$ is $R$-acyclic. The comprehensive-factorization characterization proves totality. The last equality was proved in Theorem~\ref{C:classification}; the use of integral coefficients there is essential in detecting nontrivial cyclic groups.
\end{proof}

\subsection{The exact-colimit condition, including its necessary quantifiers}

\begin{definition}[Proved: affine linearization]\label{D:aff-definition}
For a small category $A$ and a nonzero commutative ring $R$, define $\operatorname{Aff}_R(A)$ to have the objects of $A$ and the morphisms
\[
 \operatorname{Aff}_R(A)(a,b)=
 \left\{\sum_{f\in A(a,b)}r_ff\in R[A(a,b)]:
       \text{finite support},\ \sum_fr_f=1_R\right\}.
\]
Composition is the bilinear extension of composition in $A$; augmentations multiply, so composition preserves augmentation $1$. Each identity is the original identity with coefficient $1_R$. This is an ordinary small category, generally not an additive category. Its connected components are those of $A$.
\end{definition}

\begin{theorem}[Proved: the affine exact-colimit criterion]\label{D:affine-exactness}
For every small category $A$, including the empty category, and every nonzero commutative ring $R$, the following conditions are equivalent:
\begin{enumerate}
\item $\colim_A:(R\text{-}\mathrm{Mod})^A\to R\text{-}\mathrm{Mod}$ is exact.
\item Every connected component of $\operatorname{Aff}_R(A)$ is filtered.
\end{enumerate}
In explicit terms, condition (2) requires common successors within each component and an augmentation-one affine equalizer for every pair of parallel \emph{affine} arrows. It is not merely the condition in Lead~2 for one pair of original arrows. For $R=\mathbb Z$ this is the Isbell--Mitchell affine criterion \citep{IsbellMitchell1973Exact,IsbellMitchell1976Fixed}; a related finite-family formulation is verified in \citet[Section~1, unnumbered theorem]{Isbell1968Exact}. We supply the proof rather than relying on an unavailable theorem number from the original AMS papers.
\end{theorem}
\begin{proof}
Every functor $F:A\to R\text{-}\mathrm{Mod}$ extends to a functor $\bar F:\operatorname{Aff}_R(A)\to R\text{-}\mathrm{Mod}$ by
$\bar F(\sum r_ff)=\sum r_fF(f)$. This is functorial because $R$ is commutative and the maps are $R$-linear. Natural transformations extend in the same way, and pointwise exact sequences remain pointwise exact. A cocone $F\to\underline V$ respects every original arrow if and only if it respects every affine arrow: the latter follows from $\sum r_f=1_R$. Hence the identity on the individual values induces a canonical natural isomorphism
\[
 \colim_A F\cong\colim_{\operatorname{Aff}_R(A)}\bar F.
\]
Filtered colimits of modules and direct sums of modules are exact. Condition (2), followed by taking the direct sum over components, therefore implies (1).

Assume (1). We may work in one connected component of $A$, since extension by zero shows that its colimit functor is exact. For an object $a$, let $U_a$ be the set of objects receiving an arrow from $a$, and let $R_a\subseteq\underline R$ be the subfunctor taking value $R$ on $U_a$ and zero elsewhere. It is a subfunctor because $U_a$ is a cosieve. Its colimit is $R$, and the inclusion induces the identity map to $\colim_A\underline R=R$: the full subcategory $U_a$ is nonempty and connected, and all its coefficient maps are identities. Exact additive functors preserve finite limits, hence intersections of subobjects. Therefore
\[
 \colim_A(R_a\cap R_b)
 \cong(\colim_A R_a)\cap(\colim_A R_b)=R.
\]
Since $R\ne0$, the intersection functor cannot be zero, so $a,b$ have a common successor. Iteration supplies a common successor for any finite family of objects in this component.

Let $\alpha,\beta:p\to q$ be parallel affine arrows and write $\delta=\alpha-\beta\in R[A(p,q)]$, whose augmentation is zero. Let $P=R[A(p,-)]$ and let $I\subseteq P$ be the subfunctor generated, as a module-valued functor, by $\delta\in P(q)$. Thus $I(r)$ consists of finite linear combinations of $c\delta$ with $c:q\to r$. Exactness makes $\colim I\to\colim P=R$ injective. Its image is zero, because $\delta$ has augmentation zero and generates $I$. Hence $\colim I=0$.

We use the explicit direct-sum presentation of a colimit. The element $\delta$ in the summand $I(q)$ is consequently a finite sum
\[
 \delta=\sum_{j=1}^N
   \bigl(I(v_j)x_j-x_j\bigr)
 \quad\text{in }\bigoplus_{r\in A}I(r),
\]
where $v_j:s_j\to t_j$, $x_j\in I(s_j)$, and signs and scalar coefficients have been absorbed into the $x_j$. Choose a common successor $s$ of $q$ and all the $s_j,t_j$, and arrows $d_r:r\to s$ for these finitely many objects. Applying the direct-sum homomorphism whose $r$-component is $I(d_r)$ gives
\[
 d_q\delta=\sum_j(d_{t_j}v_j-d_{s_j})x_j.
\]
Write $x_j=\sum_l r_{jl}c_{jl}\delta$ with $c_{jl}:q\to s_j$. Then
\[
 t=d_q-\sum_{j,l}r_{jl}(d_{t_j}v_j-d_{s_j})c_{jl}
       \in R[A(q,s)]
\]
has augmentation $1_R$ and satisfies $t\delta=0$. Thus $t\alpha=t\beta$ in $\operatorname{Aff}_R(A)$. Combined with common successors, this is filteredness. If $A$ is empty, both conditions are vacuously satisfied: its colimit functor has the zero module as its only value.
\end{proof}

\begin{corollary}[Proved: exact module colimits imply total homology asphericity]\label{D:exact-R-sufficient}
If $A$ is nonempty and connected and $\colim_A:(R\text{-}\mathrm{Mod})^A\to R\text{-}\mathrm{Mod}$ is exact, then $A$ is totally $W_R$-aspherical. For $R=\mathbb Z$ the conclusion strengthens to total $\Winfty$-asphericity.
\end{corollary}
\begin{proof}
The direct-sum functor $\Sigma_H$ in Proposition~\ref{C:exact-sufficient} is exact for modules over $R$ as well. It makes the colimit functor on each component of each $\operatorname{El}_A(H)$ exact. Applying it to the bar resolution of the constant $R$-module proves $\mathcal P_R(A)$. Theorem~\ref{D:R-action} applies. The integral strengthening is Theorem~\ref{C:classification}.
\end{proof}

\begin{example}[Refuted: the literal homology half of Lead 4]\label{D:disconnected-counterexample}
Let $D$ have two objects $0,1$ and only their identity arrows. Its colimit functor on abelian groups is $(M_0,M_1)\mapsto M_0\oplus M_1$, which is exact. Nevertheless $D$ is not totally $W_{\mathbb Z}$-aspherical: the identity is final and $H_0(BD;\mathbb Z)=\mathbb Z^2\ne\mathbb Z$. The empty category provides another counterexample, with exact zero colimit. Thus the literal equivalence stated in Lead~4 is false; nonemptiness and connectedness cannot be suppressed. The finite category and its composition law are checked in \texttt{cd\_finite\_checks.py}.
\end{example}

\begin{remark}[Open: the corrected connected converse]\label{D:exact-converse-open}
The remaining implication is the precise statement
\[
 \left.
 \begin{array}{c}
 A\text{ small, nonempty, connected},\\
 A\text{ totally }\Winfty\text{-aspherical}
 \end{array}\right\}
 \quad\Longrightarrow\quad
 \colim_A:\Ab^A\to\Ab\text{ exact}.
\]
Equivalently, by Theorems~\ref{C:classification} and \ref{D:affine-exactness}, does $\mathcal P_\infty(A)$ force $\operatorname{Aff}_{\mathbb Z}(A)$ to be filtered? We have not proved this implication or produced a counterexample. This is an unresolved implication of the present work, not an assertion of a currently open problem in the literature. The attempted reduction through the acyclicity of $\mathbb Z[H]$ fails because these diagrams are not generally projective and their acyclicity does not supply acyclicity of arbitrary module diagrams. Lead~2 supplies affine equalizers of original parallel arrows, but exactness requires equalizers of all affine pairs. The proof in Lemma~\ref{C:monoid-lemma} bridges this gap for monoids and, by Theorem~\ref{C:weak-terminal}, for categories with a weakly terminal object. It also holds for $\Delta_{\mathrm{inj}}$ by its explicit affine cones. The first size not covered by the finite-object theorem is the case of countably infinitely many objects with no weakly terminal object; no particular smallest counterexample is claimed.

For general nonzero $R$, the corresponding unproved necessity is $\mathcal P_R(A)$, with $A$ connected and nonempty, implying filteredness of $\operatorname{Aff}_R(A)$. Neither assertion is replaced here by its sufficient direction.
\end{remark}

\subsection{Groups separate integral and other homology localizers}

\begin{proposition}[Proved: a subgroup criterion for groups]\label{D:group-criterion}
Let $G$ be a small group, regarded as a one-object category. It is totally $W_R$-aspherical if and only if every subgroup $K\leq G$ is $R$-acyclic. If $G$ is finite, this is equivalent to $|G|\in R^\times$. In particular a nontrivial finite group is totally $W_{\mathbb Q}$-aspherical but is not totally $\Winfty$-aspherical.
\end{proposition}
\begin{proof}
A covariant set-valued functor on $G$ is a left $G$-set. Its connected categories of elements are its transitive action groupoids. Choosing one object in a transitive orbit identifies its automorphism group with the corresponding stabilizer $K$, and its inclusion is an equivalence of categories. Every subgroup is realized by the coset action $G/K$. The first equivalence is therefore exactly Theorem~\ref{D:R-action}.

If $G$ is finite and $|G|$ is invertible, the element $e=|G|^{-1}\sum_{g\in G}g\in R[G]$ splits the augmentation as a homomorphism of right $R[G]$-modules: $eg=e$ and the augmentation of $e$ is $1$. Thus the trivial module $R$ is a direct summand of the free module $R[G]$. Tensoring with it, which computes coinvariants and hence $\colim_G$, is exact. Corollary~\ref{D:exact-R-sufficient} proves totality.

Conversely, if a prime $p$ divides $|G|$, Cauchy's theorem supplies a subgroup $C_p$. The standard periodic resolution of $C_p$ gives
$H_1(BC_p;R)=R/pR$. Totality forces this group to vanish, so $p$ is a unit in $R$. Every prime factor of $|G|$ is therefore a unit, and so is $|G|$. A nontrivial group has nontrivial fundamental group in its classifying space, which excludes total $\Winfty$-asphericity already for the identity final image. The order-two multiplication table and its periodic integral and rational chain complexes are checked in \texttt{cd\_finite\_checks.py}.
\end{proof}

\subsection{Truncated homotopy localizers}

\begin{definition}[Proved: truncation convention]\label{D:n-convention}
For an integer $n\geq0$, $W_n$ denotes functors inducing bijections on connected components and isomorphisms on all based homotopy groups in degrees $1,\ldots,n$, at every basepoint. Thus $W_0$ tests only components. A category is $W_n$-aspherical precisely when its classifying space is nonempty and $n$-connected. These are fundamental localizers: weak saturation and the terminal-object axiom are immediate, and the relative axiom follows by applying $n$-truncation to the two bisimplicial comparisons used in Proposition~\ref{D:homology-localizer}. The functor to $n$-types is left adjoint to inclusion \citep[Proposition~5.5.6.18]{Lurie2009HTT} and consequently preserves homotopy colimits in the category of $n$-types, so an objectwise $W_n$-equivalence induces a $W_n$-equivalence on those comparisons.
\end{definition}

\begin{lemma}[Proved: final functors and fundamental groups]\label{D:final-pi1}
An ordinary final functor $f:A\to B$ induces a bijection on components and, on each component with a chosen object $a_0$, an epimorphism
\[
 \pi_1(BA,a_0)\longrightarrow\pi_1(BB,f(a_0)).
\]
\end{lemma}
\begin{proof}
The assertion about components follows by applying finality to the colimit of the constant singleton functor. For each object $b$ in the chosen component, choose an object $(a_b,\beta_b:b\to f(a_b))$ of $b\backslash A$. At $b=f(a_0)$ choose $(a_0,\id)$. For an arrow $h:b\to c$, the two objects $(a_b,\beta_b)$ and $(a_c,\beta_ch)$ of $b\backslash A$ are joined by a zigzag because that comma category is connected. Its projection is a path from $a_b$ to $a_c$ in the groupoid completion of $A$, and its image under $f$ is the path obtained from $h$ by conjugating its endpoints with $\beta_b$ and $\beta_c$. Concatenating these paths along any based edge loop of $NB$ gives a based edge loop of $NA$ with the same image in $\pi_1(BB)$. The nerve's triangular relations are respected because each comma-category equality is an equality of composites in $B$. This proves surjectivity.
\end{proof}

\begin{theorem}[Proved: the cases $W_0$ and $W_1$]\label{D:zero-one}
A small category is totally $W_0$-aspherical if and only if it is nonempty and connected. It is totally $W_1$-aspherical if and only if its classifying space is nonempty and simply connected.
\end{theorem}
\begin{proof}
The identity final functor proves necessity in both cases. Lemma~\ref{D:final-pi1} proves sufficiency: a bijection on components preserves the unique component, and an epimorphic image of the trivial fundamental group is trivial.
\end{proof}

\begin{theorem}[Proved: the cases $W_n$, $n\geq2$]\label{D:truncated-classification}
For $n\geq2$, a small category $A$ is totally $W_n$-aspherical if and only if $BA$ is $n$-connected and every connected component of every $\operatorname{El}_A(H)$, for covariant $H:A\to\Set$, is $(n-1)$-connected.
\end{theorem}
\begin{proof}
Assume totality. The identity final image makes $BA$ $n$-connected. For a connected component of an element category, the collage gives, by reduced Mayer--Vietoris and $n$-connectedness of both $BA$ and its final image, vanishing of integral reduced homology in degrees $0,\ldots,n-1$. This truncated homological condition is inherited by categories of elements, by the explicit construction in Lemma~\ref{C:inherit-elements}. In particular it holds for every connected cyclic-subgroup covering of any such component. Since $n-1\geq1$, the argument in Lemma~\ref{C:cyclic-covers} proves that the component is simply connected. The Hurewicz theorem now makes it $(n-1)$-connected.

Conversely, if $f:A\to B$ is final, its comma categories are $(n-1)$-connected by the stated action condition. Lemma~\ref{C:homology-A}, with $d=n$, makes $H_j(Bf;\mathbb Z)$ an isomorphism for $j<n$ and an epimorphism for $j=n$. Thus $BB$ has zero reduced homology through degree $n$. It is connected and simply connected by Lemma~\ref{D:final-pi1}. Applying Hurewicz in the first possible nonzero degree shows that $BB$ is $n$-connected. This is the required totality.
\end{proof}

\begin{corollary}[Proved: truncated totality with a weakly terminal object]\label{D:truncated-finite}
For every $n\geq2$, a small category with a weakly terminal object is totally $W_n$-aspherical if and only if it is filtered. The same equivalence holds for every category with finitely many objects, including the empty category by separate inspection. It fails for $n=1$: the contractible nonfiltered right-zero monoid of Part~B is totally $W_1$-aspherical.
\end{corollary}
\begin{proof}
For $n\geq2$, Theorem~\ref{D:truncated-classification} makes every connected category of elements simply connected, so its first integral homology is zero. The proof of common successors in Lemma~\ref{C:directed-affine} needs only connectedness and vanishing first homology of $BA$ and of the doubled category of elements, not full contractibility. Proposition~\ref{B:lead2} therefore gives the same affine equalizers as before. The entire proof of Theorem~\ref{C:weak-terminal} now applies unchanged to these specified inputs. Filteredness implies total $\Winfty$-asphericity and therefore total $W_n$-asphericity. The final assertion follows from Theorem~\ref{D:zero-one} and the explicit example in Part~B.
\end{proof}

\begin{proposition}[Proved: dependence on the localizer and on the category]\label{D:general-localizer}
If $W\subseteq V$ are fundamental localizers, every totally $W$-aspherical category is totally $V$-aspherical. In particular every totally $\Winfty$-aspherical category is totally $W$-aspherical for every fundamental localizer satisfying exactly the axioms fixed in Part~A. Total $\Winfty$-asphericity is not invariant under a weak homotopy equivalence of categories.
\end{proposition}
\begin{proof}
Every tested category of elements whose map to the point lies in $W$ also has that map in $V$. The special case follows from the minimality theorem $\Winfty\subseteq W$ cited with its exact hypotheses in Part~A. For the last assertion, the nonfiltered three-element right-zero monoid of Part~B has contractible classifying space but is not total, whereas the point is filtered and total. The specified functor from that monoid to the point is a weak homotopy equivalence. This also places the same obstruction recorded by the companion at XII ms.~120: totality depends on ordinary final images and on set-valued actions, not only on the classifying-space homotopy type.
\end{proof}

\begin{remark}[Proved: the zero coefficient ring]\label{D:zero-ring}
If the zero ring is allowed, every homology group with its coefficients is zero, so $W_R$ is the class of all functors and every small category, including the empty category, is totally $W_R$-aspherical. Its module category is the zero category, and every colimit functor is exact. The nonzero hypothesis in the preceding affine and connectedness criteria is therefore necessary. The ring-theoretic statements above have been stated for commutative unital rings; no unannounced extension of their affine-linearization proof to noncommutative coefficients is being used.
\end{remark}
